% Pandoc template for the fed_fraud documents.
% Based on proposal_gen/templates/detailed_proposal_template/detailed_1_classic.tex
% (Classic report): centered title page, table of contents, ruled section
% headings and the Geodesic house footer. Fonts are Latin Modern (lmodern).
% Two layouts: the full report (title page and contents) and, with
% -V brief, a compact one-page-brief header in the same style.
\documentclass[letterpaper,11pt]{article}
\usepackage[T1]{fontenc}
\usepackage[utf8]{inputenc}
\usepackage{lmodern}
\usepackage[letterpaper,top=0.9in,bottom=1.2in,left=1in,right=1in,footskip=0.6in]{geometry}
\usepackage[table]{xcolor}
\usepackage{geodesicbrand}
\usepackage{amsmath,amssymb}
\usepackage{enumitem}
\usepackage{array}
\usepackage{calc}
\usepackage{longtable}
\usepackage{makecell}
\usepackage[most]{tcolorbox}
\usepackage{titlesec}
\usepackage{needspace}
\usepackage{etoolbox}
\usepackage{hyperref}
\hypersetup{colorlinks=true,urlcolor=coolblack,linkcolor=coolblack,filecolor=geoblue,
  pdftitle={Red-team playbook},pdfauthor={\geofirm}}
\urlstyle{same}

\setlist{itemsep=2pt,leftmargin=1.4em}
\setlength{\parindent}{0pt}
\setlength{\parskip}{6pt}
\providecommand{\tightlist}{\setlength{\itemsep}{0pt}\setlength{\parskip}{0pt}}
\setcounter{secnumdepth}{0}

% Section style: house rule under each heading
\titleformat{\section}{\Large\bfseries\color{coolblack}}{\thesection}{0.6em}{}[{\color{coolblack}\titlerule[0.5pt]}]
\titleformat{\subsection}{\large\bfseries\color{geomid}}{\thesubsection}{0.6em}{}
\titleformat{\subsubsection}{\normalsize\bfseries\color{geomid}}{\thesubsubsection}{0.6em}{}
\titlespacing*{\section}{0pt}{10pt}{6pt}

% Tables are written by latex/filters.lua (navy header, banded rows)
\setlength{\LTpre}{6pt}\setlength{\LTpost}{6pt}

\begin{document}

% BRIEF HEADER
\thispagestyle{geodesic}
\begin{minipage}[c]{0.13\textwidth}\includegraphics[width=\linewidth]{icon.png}\end{minipage}\hfill
\begin{minipage}[c]{0.84\textwidth}
  {\textsc{\large \geofirm}}\par\vspace{2pt}
  {\small\color{geogray}Federal fraud expansion\enspace\textperiodcentered\enspace When: \compiletime}
\end{minipage}\par\vspace{6pt}
{\color{coolblack}\rule{\linewidth}{0.3mm}}\par\vspace{2pt}
{\LARGE\bfseries\color{coolblack} Red-team playbook}\par
\vspace{2pt}{\large\color{geomid} How fraudsters adapt, and how the controls hold}\par
\vspace{-2pt}{\color{coolblack}\rule{\linewidth}{0.3mm}}\par\vspace{4pt}
\pagestyle{geodesic}

\textbf{Status:} Draft v0.1 (2026-11-04), for Steve's review. Extension
10 of 10 (see the \href{../extensions.pdf}{extensions overview}). This
playbook is about the durability of our detection controls. It describes
adaptations at the level an agency risk office would discuss; it is not
a guide to committing fraud.

\hypertarget{bottom-line}{%
\section*{Bottom line}\label{bottom-line}}
\addcontentsline{toc}{section}{Bottom line}

Every published control teaches fraudsters what to avoid. A proposal
that ignores this will be challenged by the first program official who
has watched a scheme move from one program to the next. This playbook
takes each control in the plan, names the obvious adaptation, says what
signal the adaptation itself leaves, and states the counter. Four
principles come out of it:

\begin{enumerate}
\def\labelenumi{\arabic{enumi}.}
\tightlist
\item
  \textbf{Layer signals that are costly to fake together.} Avoiding one
  signal (a shared address) usually creates another (many new addresses
  registered the same week). Leads need two or more independent signals.
\item
  \textbf{Watch for the adaptation, not just the scheme.} Displacement
  is itself a signal: when a control tightens in one program, look for
  the same actors in the next program and the next state.
\item
  \textbf{Keep thresholds and weights private.} Publish methods at the
  level of the handbook, never the exact cut-offs or the feature
  weights.
\item
  \textbf{Re-validate on a schedule.} Every screen is re-tested
  quarterly on new confirmed cases, and drift in hit rates triggers
  review.
\end{enumerate}

\hypertarget{adaptations-by-control}{%
\section*{1. Adaptations by control}\label{adaptations-by-control}}
\addcontentsline{toc}{section}{1. Adaptations by control}

\begingroup\small\setlength{\tabcolsep}{4pt}\renewcommand{\arraystretch}{1.25}
\rowcolors{2}{white}{geolight}
\begin{longtable}[l]{@{}>{\raggedright\arraybackslash}p{\dimexpr 0.189\linewidth-2\tabcolsep}>{\raggedright\arraybackslash}p{\dimexpr 0.270\linewidth-2\tabcolsep}>{\raggedright\arraybackslash}p{\dimexpr 0.270\linewidth-2\tabcolsep}>{\raggedright\arraybackslash}p{\dimexpr 0.270\linewidth-2\tabcolsep}@{}}
\rowcolor{coolblack}\leavevmode\color{white}\bfseries Control & \leavevmode\color{white}\bfseries Likely adaptation & \leavevmode\color{white}\bfseries Trace the adaptation leaves & \leavevmode\color{white}\bfseries Counter \\
\endhead
Shared-address and shared-officer networks & One address and one nominee officer per entity & Bursts of new registrations, virtual-office and mailbox addresses, nominees appearing across many unrelated sectors & Co-location test on address type; registration-timing clusters; nominee screens across registries \\
Bank-account sharing & One account per entity, money moved on afterward & Accounts opened days before first payment, at the same few institutions; rapid outbound transfers & Account-age and account-owner verification at enrollment (Phase 2, with Treasury) \\
Exclusion matching on name and identifiers & New entity with a relative or associate as owner & Same addresses, phones, staff or patients as the excluded entity & Linked-identity matching across attributes, not names; network seed propagation \\
Death and identity checks & Synthetic identities (real SSN fragments plus invented details) & Thin credit histories, identities that appear in only one program, shared devices & Cross-program identity age; device and contact sharing (Phase 2) \\
Billing-ramp and peer-outlier screens & Grow slowly; stay just inside peer ranges; spread billing across many entities & Many entities each at the 85th to 95th percentile; identical billing mixes & Distribution-shape tests across a group; cluster-level scoring rather than entity-level \\
Threshold bunching tests & Vary amounts randomly below the threshold & Too-uniform distributions; digit tests on the last two digits & Second-order and last-digit tests; bunching on a wider window \\
Benford and round-number tests & Use invented ``random-looking'' amounts & People are poor at random: too few repeats, too even digits & Number duplication and summation tests; compare to the entity's own history \\
Sealed-bid screens & Fake competition with a wider spread of cover bids & Cover bidders that never win, bid only against each other, or share preparers & Win-loss networks across tenders; rotation tests over time \\
Set-aside size and pass-through checks & Hide the large partner as a ``mentor'' or through a joint venture & Mentor-protege pairs where the protege does little of the work; sudden revenue jumps & Subaward and payroll share checks; front-company profile \\
Cross-agency matching & Move to a program or state not yet in the match & Same actors appearing in new programs after a control goes live & Displacement watch list; prioritize the next programs for matching \\
Tip triage & Flood hotlines with low-quality tips about competitors or to bury a real tip & Bursts of similar tips from one source or template & Duplicate and campaign detection; tip source analysis \\
Pre-payment checks & Target programs or periods with pay-first rules (emergencies) & Spikes in new payees in emergency programs & Emergency playbooks with identity verification that is fast, not absent \\
\end{longtable}\endgroup

\hypertarget{pushback-we-should-expect-and-the-answer}{%
\section*{2. Pushback we should expect, and the
answer}\label{pushback-we-should-expect-and-the-answer}}
\addcontentsline{toc}{section}{2. Pushback we should expect, and the
answer}

\begingroup\small\setlength{\tabcolsep}{4pt}\renewcommand{\arraystretch}{1.25}
\rowcolors{2}{white}{geolight}
\begin{longtable}[l]{@{}>{\raggedright\arraybackslash}p{\dimexpr 0.500\linewidth-2\tabcolsep}>{\raggedright\arraybackslash}p{\dimexpr 0.500\linewidth-2\tabcolsep}@{}}
\rowcolor{coolblack}\leavevmode\color{white}\bfseries Pushback & \leavevmode\color{white}\bfseries Answer \\
\endhead
``Fraudsters will just adapt.'' & Yes; that is why controls are layered and re-validated, and why adaptation leaves its own trace (table above). Every adaptation raises the cost of the fraud. \\
``Your false positives will hurt legitimate providers and beneficiaries.'' & No action is taken on a score. Leads go to trained reviewers, each screen states its false-positive rate, and release requires 5x lift and bias checks. \\
``Publishing your methods teaches them.'' & We publish the method class, not thresholds or weights. Academic and GAO literature already describes the methods; the advantage is in data and execution. \\
``We already do this.'' & Then the pilot measures our marginal value directly: the new-detection rate against what the agency already knows. If it is below break-even, we stop. \\
``The data will never be shared.'' & Most of Phase 1 runs on public data, and the legal routes for Phase 2 exist (see the \href{../policy/legal_track.pdf}{policy and legal track}). \\
\end{longtable}\endgroup

\hypertarget{running-a-red-team-exercise}{%
\section*{3. Running a red-team
exercise}\label{running-a-red-team-exercise}}
\addcontentsline{toc}{section}{3. Running a red-team exercise}

Twice a year, and before any screen's first release:

\begin{enumerate}
\def\labelenumi{\arabic{enumi}.}
\tightlist
\item
  A team member not involved in building the screen takes the role of a
  fraud organizer and writes down how they would avoid it, using only
  what is public.
\item
  The team builds a synthetic version of the adapted scheme and runs the
  full screen set on it.
\item
  Each adaptation is scored: caught by another screen, caught by the
  adaptation's own trace, or missed.
\item
  Misses become backlog items with an owner; results go into the
  screen's model documentation.
\end{enumerate}

\hypertarget{questions-for-steve}{%
\section*{4. Questions for Steve}\label{questions-for-steve}}
\addcontentsline{toc}{section}{4. Questions for Steve}

\begin{enumerate}
\def\labelenumi{\arabic{enumi}.}
\tightlist
\item
  Should we invite an OIG investigator or a former prosecutor to the
  first red-team session?
\item
  How much of the method do we put on the public parts of the site? The
  current site is password-protected; this playbook assumes the
  thresholds stay off any public page.
\end{enumerate}

\end{document}
